% !TeX root = ../Book4.tex
\begin{figure}[htbp]
	\centering
	\begin{tikzpicture}[font=\small,x=1.15cm,y=1.05cm]
		\fill[GeometryBlueLight] (-1.2,0) -- (-1.2,2.2) arc[start angle=180,end angle=360,x radius=1.2,y radius=0.32] -- (1.2,0) arc[start angle=0,end angle=-180,x radius=1.2,y radius=0.32];
		\draw[GeometryGray,dashed] (1.2,0) arc[start angle=0,end angle=180,x radius=1.2,y radius=0.32];
		\draw[GeometryBlue,thick,->] (-1.2,0) arc[start angle=180,end angle=355,x radius=1.2,y radius=0.32];
		\draw[GeometryBlue,thick] (0,2.2) ellipse[x radius=1.2,y radius=0.32];
		\draw[GeometryGray] (-1.2,0)--(-1.2,2.2) (1.2,0)--(1.2,2.2);
		\draw[GeometryBlue,thick,->] (0,-0.32)--(0,1.88);
		\node[right] at (0,1.1) {\(t\)};
		\node[below] at (0,-0.4) {\(s\in\mathbb R/\mathbb Z\)};
		\node[left] at (-1.2,0) {\(\gamma\)};
		\node[left] at (-1.2,2.2) {\(\phi_1\circ\gamma\)};
		\node[right,align=left] at (1.65,1.1) {\(F(t,s)=\phi_t(\gamma(s))\)\\[3pt]\(\displaystyle\int F^*\omega=\langle\operatorname{Flux},[\gamma]\rangle\)};
	\end{tikzpicture}
	\caption[闭路运动与通量]{参数圆柱的横向是闭路参数，纵向是同位时间。通量在该闭路上的取值等于按 \(\mathrm{d}t\wedge\mathrm{d}s\) 定向的拉回辛面积；图中的圆柱表示参数域，其在流形中的像可以自交。}
	\label{g4:fig:0903:swept-cylinder}
\end{figure}
